\documentclass[12pt]{article}
\usepackage[a4paper, margin=2.5cm]{geometry}
\usepackage{hyperref}
\usepackage{booktabs}
\usepackage{array}
\usepackage{parskip}

\begin{document}

\noindent
\textbf{To:}\\
The Editor\\
Journal of Physics Communications\\
IOP Publishing\\
Temple Circus, Temple Way\\
Bristol BS1 6HG, United Kingdom

\bigskip
\noindent August 2026

\bigskip
Dear Editor,

\section*{Submission Statement}

I hereby submit the manuscript

\begin{center}
\textit{``Resonance Field Theory: Axiomatics, Fundamental Formula, and Empirical Validation''}
\end{center}

\noindent by Dominic-Ren\'e Schu (ORCID: 0009-0004-9769-9061) for consideration as an
\textbf{Original Research Article} in the \textit{Journal of Physics Communications}.

\section*{Core Statement}

Resonance Field Theory (RFT) is an axiomatic framework that derives a phase-dependent
extension of the Planck relation,
$E = \pi \cdot \varepsilon(\Delta\varphi) \cdot \hbar \cdot f$,
from seven hierarchically structured axioms. The central result is that the factor $\pi$
in this formula is not a free normalization parameter but emerges geometrically from the
stationary-phase contribution of the action integral over a half-cycle of phase space
(RT-01, RT-01b); and that the coupling efficiency
$\varepsilon(\Delta\varphi) = \cos^2(\Delta\varphi/2)$ is the unique function satisfying
all boundary conditions within the $k\!=\!1$ irreducible representation of
$\mathrm{U}(1) \subset G_{\mathrm{sync}}$, the theory's symmetry group (RT-02). The
theory is empirically validated across four independent physical domains --- particle
physics, cosmology, nuclear physics, and quantum mechanics --- and two concrete
falsifiable experimental proposals are presented with numerically specified decision
criteria.

\textit{Journal of Physics Communications} is the appropriate venue because the
manuscript communicates a significant original result that simultaneously (a) introduces
a formally structured axiomatic framework, (b) presents cross-domain empirical
validation with quantified significance, and (c) provides experimentally testable
predictions accessible to existing laboratory infrastructure.

\section*{Distinction from Prior Work}

RFT is not a reformulation of existing oscillator physics. Three features distinguish
it from prior frameworks:

\textbf{$\pi$ as a derived geometric constant (RT-01, RT-01b).}
In the standard Planck relation $E = \hbar\omega$, the factor $2\pi$ is built into the
definition of $\omega$. In RFT, $\pi$ appears explicitly as a result of integrating the
saddle-point contribution of the action integral
$S[\varphi] = \int_0^{\pi} \varepsilon(\varphi)\,d\varphi$ over one half-period of the
resonance cycle. The numerical confirmation ($|c_3 + c_4| \approx 5.5 \times 10^{-11}$,
RT-01b) shows this is not an arbitrary normalization. The falsification condition is
explicit: any measurement demonstrating that the phase-integrated coupling energy
deviates from $\pi \cdot \varepsilon \cdot \hbar \cdot f$ beyond the numerical precision
of the saddle-point calculation would refute the derivation.

\textbf{$\varepsilon = \cos^2(\Delta\varphi/2)$ as representation-theoretically forced
(RT-02).}
Within the symmetry group
$G_{\mathrm{sync}} \cong \mathbb{R}^+_\times \times \mathrm{U}(1) \times
\mathrm{Aff}^+(\mathbb{R})$, the $k\!=\!1$ irreducible representation of $\mathrm{U}(1)$
admits exactly one function $F\colon[0,2\pi]\to[0,1]$ satisfying $\varepsilon(0)=1$,
$\varepsilon(\pi)=0$, parity $\varepsilon(\Delta\varphi)=\varepsilon(-\Delta\varphi)$,
and non-mixing of higher representations ($k \geq 2$). That function is
$\cos^2(\Delta\varphi/2)$. The uniqueness theorem is proved in \S3.3 of the manuscript.

\textbf{A3 as a corollary, not an independent axiom (RT-02, RT-35).}
The resonance condition $|f_1/f_2 - m/n| < \delta$ is derived as a corollary of
Axiom~A7 (invariance under $G_{\mathrm{sync}}$) and is no longer counted as an
independent postulate. This reduces the axiom count from seven to six effective
independent postulates.

\textbf{Falsifiable predictions with quantified criteria.}
Two experimental proposals include specific falsification thresholds: (I) SNR $\geq 3\sigma$
for Am-241 photo-excitation at ELI-NP with 100~h beam time
(SNR$_{\mathrm{median}} \approx 10\sigma$); (II) center-of-mass shift
$|\Delta\langle x\rangle| \approx 2.0\cdot\lambda\;\mu\mathrm{m}$ for ${}^{87}$Rb BEC
in a harmonic trap at $\Delta\varphi = \pi$, testable with existing BEC apparatus.

\section*{Empirical Reach}

\begin{table}[h!]
\centering
\begin{tabular}{p{3cm} p{5cm} p{5.5cm}}
\toprule
Domain & Method & Key result \\
\midrule
Particle physics & 1,500,000 MC sim.\ on CMS Open Data dielectron events &
  5 resonances detected (empirical $p = 0$, stable across 3 bandwidths, 10~seeds) \\
Cosmology & 1,530 coupled FLRW simulations &
  $\eta \approx \cos^2(\Delta\varphi/2)$ emergent; $\Delta d_\eta > 6\sigma$;
  $\Delta\chi^2 = +16$ vs.\ Planck-2018 CMB \\
Nuclear physics & Resonance reactor, $\kappa = 1$ (no free parameter) &
  Fission break-even $Q \approx 1.0$; $\lambda_{\mathrm{eff}}/\lambda_0 = 7{,}872$ \\
Quantum mechanics & Schr\"odinger simulation (5 derivation steps) &
  Fidelity $= 1.0$ (all 4 scenarios); $1-F \sim \lambda^2$ confirmed \\
\bottomrule
\end{tabular}
\end{table}

\section*{Open Science Statement}

All derivations, source code, simulation scripts, and raw data are openly available at:
\begin{center}
\url{https://github.com/DominicReneSchu/RFT}
\end{center}
\noindent This includes all RT-series analyses (RT-01 through RT-38), the double pendulum
experiment protocol (RT-38, reproducible with approximately 200~EUR and a smartphone),
and the full backtest infrastructure. The repository is public and version-controlled.

\section*{Suggested Reviewers}

\noindent\textbf{Particle Physics / Resonance Physics:}

\begin{itemize}
\item Dr.\ G\"unter Dissertori, ETH Z\"urich (CMS collaboration, resonance searches in
  dielectron/dimuon channels): \texttt{dissertori@phys.ethz.ch}
\item Dr.\ Markus Klute, MIT (Higgs boson physics, resonance structure in CMS data):
  \texttt{mklute@mit.edu}
\end{itemize}

\noindent\textbf{Mathematical Physics (Group Theory / Representation Theory):}

\begin{itemize}
\item Prof.\ Gerd Rudolph, Universit\"at Leipzig (gauge field theory,
  group-theoretic structures): \texttt{rudolph@itp.uni-leipzig.de}
\end{itemize}

\noindent\textbf{Cosmology (FLRW / CMB):}

\begin{itemize}
\item Dr.\ Licia Verde, Universitat de Barcelona / ICREA (CMB power spectrum analysis,
  Hubble tension): \texttt{liciaverde@icc.ub.edu}
\end{itemize}

\noindent\textbf{Quantum Mechanics / Atomic Physics (BEC / Interferometry):}

\begin{itemize}
\item Prof.\ Markus Oberthaler, Universit\"at Heidelberg (BEC physics,
  matter-wave interferometry with ${}^{87}$Rb): \texttt{markus.oberthaler@kip.uni-heidelberg.de}
\end{itemize}

\section*{Excluded Reviewers}

No reviewers are formally excluded. There are no personal, commercial, or intellectual
conflicts that would warrant exclusion.

\section*{Author Information}

\begin{tabular}{ll}
\textbf{Author:} & Dominic-Ren\'e Schu \\
\textbf{Affiliation:} & Independent Researcher \\
\textbf{Email:} & dominic.rene.schu@gmail.com \\
\textbf{ORCID:} & 0009-0004-9769-9061 \\
\textbf{Co-authors:} & None \\
\textbf{Conflicts of interest:} & None \\
\textbf{Funding:} & No external funding \\
\textbf{Data availability:} & \url{https://github.com/DominicReneSchu/RFT} \\
\end{tabular}

\bigskip
\noindent I confirm that this manuscript is original, has not been published elsewhere,
and is not currently under consideration at any other journal.

\bigskip\bigskip
\noindent Yours sincerely,

\bigskip\bigskip
\noindent Dominic-Ren\'e Schu\\
August 2026

\end{document}
